\section{Manajemen Tabel Simbol dan Penanganan Kesalahan}
\label{sec:ch03-simbol}

\subsection{Tabel Simbol}

\begin{definisi}[Tabel simbol]
Tabel simbol adalah struktur data yang memetakan identifier ke
informasi tentang identifier tersebut, seperti tipe, scope, alamat
memori, dan kategori (variabel, fungsi, parameter).
\end{definisi}

Tabel simbol mendukung operasi \code{insert}, \code{lookup}, dan
\code{delete}. Pada bahasa dengan scope bersarang, tabel dapat
diimplementasikan sebagai stack tabel atau satu tabel dengan rantai
scope. Informasi yang disimpan biasanya mencakup nama, tipe, ukuran,
alamat, dan referensi ke deklarasi.

Tabel simbol dipakai oleh hampir semua fase. Parser menggunakannya
untuk membedakan identifier dan kata kunci; analisis semantik
menggunakannya untuk memeriksa deklarasi dan tipe; pembangkitan kode
menggunakannya untuk menentukan alamat variabel. Karena itu, tabel
simbol sering disebut sebagai pusat informasi kompilator.

\subsection{Jenis Kesalahan}

Kesalahan program dapat diklasifikasikan menurut fase yang
mendeteksinya:
\begin{itemize}[leftmargin=*]
  \item \textbf{kesalahan leksikal}: karakter ilegal, string tidak
        ditutup, komentar tidak ditutup;
  \item \textbf{kesalahan sintaksis}: token tidak sesuai tata bahasa,
        kurung tidak seimbang, titik koma hilang;
  \item \textbf{kesalahan semantik}: identifier tidak dideklarasikan,
        tipe tidak cocok, argumen fungsi salah.
\end{itemize}

Kesalahan \eng{run-time} seperti pembagian nol atau indeks di luar
batas bukan tanggung jawab \eng{front-end}; kesalahan itu muncul saat
program dijalankan, bukan saat program diterjemahkan.

Compiler yang baik tidak berhenti pada kesalahan pertama. Ia
melaporkan beberapa kesalahan, berusaha pulih, dan melanjutkan
analisis. Strategi pemulihan yang umum adalah \eng{panic mode},
\eng{phrase-level recovery}, dan \eng{error production}.

\begin{table}[H]
\centering
\caption{Contoh pesan kesalahan yang informatif}
\label{tab:ch03-error}
\begin{tabularx}{\textwidth}{@{}lXX@{}}
\toprule
\textbf{Jenis} & \textbf{Pesan buruk} & \textbf{Pesan informatif} \\
\midrule
Leksikal & \code{error} & \code{baris 4, kolom 12: karakter '@' tidak dikenal} \\
Sintaksis & \code{syntax error} & \code{baris 7: diharapkan ';' setelah ekspresi} \\
Semantik & \code{type error} & \code{baris 9: tidak dapat menetapkan float ke int tanpa cast} \\
\bottomrule
\end{tabularx}
\end{table}

Pesan yang baik memuat lokasi (baris dan kolom), penyebab, dan saran
perbaikan bila memungkinkan. Informasi lokasi sangat membantu karena
mahasiswa dan pengguna dapat langsung menuju sumber masalah.

\subsection{Contoh Alur Tabel Simbol}

Perhatikan potongan program:
\begin{lstlisting}[language=C]
int x = 1;
{
    int x = 2;
    print(x);
}
print(x);
\end{lstlisting}
Saat memasuki blok, tabel simbol baru didorong ke stack. Deklarasi
\code{int x = 2} berada pada scope dalam dan menutupi (\eng{shadow})
\code{x} luar. Saat blok selesai, tabel scope dalam dibuang, sehingga
\code{print(x)} terakhir merujuk ke \code{x} luar. Perilaku ini akan
diimplementasikan pada Bab XV.
